\section{Background results}
\label{ch:required-results}

This chapter collects theorems that the proof relies on but whose proofs appear in
other papers. Each entry gives a full statement, a pointer to the literature, and
comments on its role and on formalization. The project may initially assume these as
axioms; discharging them are independent sub-projects. Three of these results are
larger sub-projects and are developed in their own sections: the computability
results in Section~\ref{ch:computability} and parallel repetition in
Section~\ref{ch:parallel-repetition}; they are included in the summary table below
for completeness. Statements are given in the
form closest to how they will be used; where the source states a bipartite version and
the project needs a synchronous one (see Section~\ref{sec:departures}), the translation
is part of the formalization task and is noted in the comments.

Each entry carries an estimate of the formalization effort:
\effortEasy{} for self-contained arguments of at most a few pages,
\effortMedium{} for substantial but well-delimited developments, and
\effortHard{} for results whose formalization is a project in itself;
\effortMathlib{} marks statements essentially available in Mathlib. The estimates
are for the statement as it will be used (including any synchronous translation) and
are of course provisional. Independently of the effort estimate, the \qlib{} badge
(introduced in Section~\ref{ch:foundations}) marks results of general
quantum-information interest that are intended to be upstreamed to QuantumLib.

\begin{center}
\small
\begin{tabular}{llll}
\hline
Result & Statement & Source & Effort \\
\hline
Gowers--Hatami & Thm.~\ref{thm:gowers-hatami} & \cite{GH17,dlS21} & \checkmark \\
Almost-synchronous correlations & Thm.~\ref{thm:almost-sync} & \cite{Vid22} & hard \\
Orthonormalization & Thm.~\ref{thm:orthonormalization} & \cite{dlS21} & \checkmark \\
Magic Square rigidity & Thm.~\ref{thm:ms-rigidity} & \cite{WBMS16} & easy \\
Low individual degree test (classical) & Thm.~\ref{thm:lidt-soundness} & \cite{JNVWY21qld} & \checkmark \\
Low individual degree test (Pauli) & Thm.~\ref{thm:qld} & \cite{JNVWY21qld} & hard \\
Direct parallel repetition (commuting-operator) & Thm.~\ref{thm:direct-repetition-co} & \cite{Vid26} & \checkmark \\
Direct parallel repetition (entangled) & Thm.~\ref{thm:direct-repetition-q} & \cite{OpenAI26} & \checkmark \\
Tracial density & Thm.~\ref{thm:tracial-density} & \cite{Lin23,Vid26} & \checkmark \\
Recursive compression lemma & Lem.~\ref{lem:recursive-compression} & \cite{MNY} & \checkmark \\
Compressibility criterion & Lem.~\ref{lem:compressible-criterion} & \cite{Lin25} & \checkmark \\
Efficient computability toolkit & Lems.~\ref{lem:universal-tm}--\ref{lem:kleene} & \cite{MNY} & hard \\
Halting undecidability & Thm.~\ref{thm:halting-undecidable} & Mathlib & \checkmark \\
Succinct Cook--Levin & Thm.~\ref{thm:succinct-sat} & \cite{Coo71,PY86} & hard \\
Self-dual normal bases & Lem.~\ref{lem:self-dual-basis} & \cite{Sho90,Len91} & medium \\
Schwartz--Zippel & Lem.~\ref{lem:schwartz-zippel} & \cite{Sch80,Zip79} & \checkmark \\
Value approximable from below & Lem.~\ref{lem:value-lower-approx} & folklore & medium \\
\hline
\end{tabular}
\end{center}

\noindent The anchored parallel repetition theorem of~\cite{BVY17}, and the entropy and
fidelity toolkit its proof needs, are not on the path to the main theorem and are not
stated in this blueprint (Remark~\ref{rem:direct-vs-anchored}).

\subsection{The Gowers--Hatami theorem}
\label{sec:rr-gowers-hatami}

\qlib

\begin{definition}[Approximate representation]
\label{def:approx-rep}
\uses{def:hs-norm}
\lean{MIPRE.IsApproxRepresentation}
\leanok
\qlib{}
Let $G$ be a finite group and $\eps \in \R$. An \emph{$\eps$-approximate
representation} of $G$ is a function $f : G \to U(d)$ such that
\begin{equation*}
  \E_{x, y \sim G} \, \mathrm{Re}\, \langle f(x) f(y), f(xy) \rangle_{hs}
  \geq 1 - \eps\;,
\end{equation*}
where $\langle \cdot, \cdot \rangle_{hs}$ is the normalized Hilbert--Schmidt
inner product of Definition~\ref{def:hs-norm}. Since each value of $f$ is
unitary, $\norm{f(x) f(y) - f(xy)}_{hs}^2 = 2 - 2\, \mathrm{Re}\, \langle f(x)
f(y), f(xy) \rangle_{hs}$, so the condition is equivalent to the defect form
$\E_{x, y \sim G} \norm{f(x) f(y) - f(xy)}_{hs}^2 \leq 2 \eps$.
\end{definition}

\begin{theorem}[Gowers--Hatami]
\label{thm:gowers-hatami}
\uses{def:hs-norm,def:approx-rep}
\lean{MIPRE.gowers_hatami}
\leanok
Let $G$ be a finite group and $f : G \to U(d)$ an $\eps$-approximate
representation of $G$. Then there exist $d' \geq d$, an isometry
$V : \C^d \to \C^{d'}$, and a representation $\rho : G \to U(d')$ such that
\begin{equation*}
  \E_{x \sim G} \, \norm{f(x) - V^\dagger \rho(x) V}_{hs}^2 \leq 2 \eps\;.
\end{equation*}
\end{theorem}

\begin{proof}
\leanok
Let $L(G, \C^d) \cong \C^{|G| d}$ carry the inner product normalized by $1/|G|$,
let $V : \C^d \to L(G, \C^d)$ be the map $u \mapsto (x \mapsto f(x) u)$, and let
$\rho$ be the right regular representation, $(\rho(x) F)(y) = F(yx)$. Then $V$ is
an isometry, and the compression $V^\dagger \rho(x) V = \E_{y \sim G}
f(y)^\dagger f(yx)$ is an average of unitaries, hence has squared
Hilbert--Schmidt norm at most $1$. Moreover
$\E_{x} \mathrm{Re}\, \langle f(x), V^\dagger \rho(x) V \rangle_{hs} =
\E_{x, y} \mathrm{Re}\, \langle f(x) f(y), f(xy) \rangle_{hs} \geq 1 - \eps$ by
hypothesis, and the claim follows by expanding
$\norm{f(x) - V^\dagger \rho(x) V}_{hs}^2$. The construction gives
$d' = |G|\, d$.
\end{proof}

\paragraph{Source.} Gowers and Hatami~\cite{GH17}. A quantitative form suited to
applications in nonlocal games, and a version adapted to the tracial setting, is given
by de la Salle~\cite{dlS21} and in~\cite{Vid22}.

\paragraph{Comments.} This is the stability result underlying all rigidity
(self-testing) arguments in the proof: an approximate representation of the Pauli group
$\F_2^{2k}$ (with its cocycle) is close to an exact one, which forces the presence of
EPR pairs. The statement above is fully formalized (statement and proof) in
\texttt{MIPRE/Background/GowersHatami/Basic.lean}, following the regular-representation
argument sketched in the proof; the intermediate statement, with the exact
representation living on the concrete index set $G \times \{1, \dots, d\}$, is
\texttt{MIPRE.gowers\_hatami\_prod}. Two remarks on the statement. (i) It asserts
only $d' \geq d$, with witness $d' = |G| d$; it does not include the dimension bound
$d' \leq (1 + C\eps) d$ of~\cite{GH17}, whose proof goes through the Fourier-analytic
route via irreducible representations. Whether the sharper bound is needed downstream
(in the Pauli-group application feeding Theorem~\ref{thm:qld}) is to be revisited; if
so it will be added as a separate strengthening. (ii) In the Lean encoding the
isometry is stored as the $d \times d'$ matrix of its adjoint, so the isometry
condition reads $V V^\dagger = \Id$ and the compression is written
$V \rho(x) V^\dagger$. The normalized Hilbert--Schmidt norm is the tracial form
needed for the synchronous framework~\cite{dlS21,Vid22}.

\subsection{Almost-synchronous correlations}
\label{sec:rr-almost-sync}

\effortHard\ \qlib

\begin{theorem}[Almost-synchronous strategies are near synchronous strategies]
\label{thm:almost-sync}
\ledgernode{1.2.1.7.2, 1.2.1.7.2.1, 1.2.1.7.2.2, 1.2.1.7.2.3}
\uses{def:sync-game,def:sync-value,def:q-value,def:inconsistency}
There are universal constants $C > 0$ and $0 < c \leq 1$ such that the following holds.
Let $\game$ be a synchronous game whose question distribution gives the diagonal
pointwise proportional weight: for some $\kappa \in (0, 1]$,
\begin{equation*}
  \mu(x,x) \;\geq\; \kappa \sum_{y \in \mX} \mu(x,y) \qquad \text{for all } x \in \mX\;.
\end{equation*}
Then, writing $\eps = 1 - \valstar(\game)$,
\begin{equation*}
  1 - \synval(\game) \;\leq\; C \bigl( \eps / \kappa \bigr)^{c}\;.
\end{equation*}
More precisely, every tensor-product strategy for $\game$ with value $1 - \eps$ is,
up to local isometries, well-approximated on average by a convex combination of
synchronous strategies each of value at least $1 - C(\eps/\kappa)^{c}$.
\end{theorem}

\paragraph{Source.} Vidick~\cite{Vid22}, building on~\cite{KPS18,PSSTW16}; the
statement above is the corollary needed here, with constants independent of the game's
size. What the source controls is not the game's value directly but the
\emph{synchronicity defect} of a strategy --- its inconsistency
(Definition~\ref{def:inconsistency}) at equal questions, taken against the marginal of
$\mu$ on $\mX$; the diagonal-weight hypothesis is what turns the game's rejection
probability $\eps$ into the bound $\eps/\kappa$ on that defect, the pairs $(x,x)$ carrying
at least a $\kappa$ fraction of the weight of $x$. Without it the two values are unrelated
(Remark~\ref{rem:almost-sync-hypothesis}).

\paragraph{Comments.} The proof uses the orthonormalization
lemma (Section~\ref{sec:rr-orthonormalization}) and a rounding argument. Used
contrapositively, as Theorem~\ref{thm:parallel-repetition} uses it, the theorem turns a
bound $\synval(\game) \leq 1 - \delta$ into $\valstar(\game) \leq 1 - \Omega(\kappa
\delta^{1/c})$: the diagonal weight multiplies the gap recovered, so a small $\kappa$ is a
loss. The hypothesis is also not preserved by repetition, which is why the soundness of
the pipeline is carried in $\valstar$ from Theorem~\ref{thm:parallel-repetition} onwards.

\begin{remark}[The diagonal-weight hypothesis cannot be dropped]
\label{rem:almost-sync-hypothesis}
\uses{thm:almost-sync,def:sync-game,def:sync-value,def:q-value}
Some lower bound on the diagonal weight is necessary, and its absence is not merely
lossy: the two values become unrelated. Let $\mX = \mA = \{0,1\}$, let $\mu$ give weight
$\frac12$ to each of $(0,1)$ and $(1,0)$ and none to the diagonal, and let $D$ accept
exactly the answer pair $(a,b) = (0,1)$ at those two question pairs and reject unequal
answers at the two diagonal pairs, as Definition~\ref{def:sync-game} requires --- which
costs nothing here, the diagonal carrying no weight. Alice answering $0$ and Bob answering
$1$, deterministically, is a tensor-product strategy of value $1$, so
$\valstar(\game) = 1$. A synchronous strategy, by contrast, is a single family
$\{M^x_a\}$ of projective measurements against a tracial state $\tau$; writing
$P = M^0_0$ and $Q = M^1_0$, its value is
\begin{equation*}
  \tfrac12 \tau\bigl( P (\Id - Q) \bigr) + \tfrac12 \tau\bigl( Q (\Id - P) \bigr)
  \;=\; \tfrac12 \bigl( \tau(P) + \tau(Q) \bigr) - \tau(PQ)
\end{equation*}
by traciality. Now $\tau(PQ) \geq \max\{0,\ \tau(P) + \tau(Q) - 1\}$: the first bound
because $\tau(PQ) = \tau(QPQ) \geq 0$, the second because
$\tau\bigl( (\Id - P)(\Id - Q) \bigr) = \tau\bigl( (\Id - Q)(\Id - P)(\Id - Q) \bigr)
\geq 0$. With $s = \tau(P) + \tau(Q)$ the value is therefore at most
$\min\{ s/2,\ 1 - s/2 \} \leq \frac12$, attained in dimension $1$ at $P = \Id$, $Q = 0$.
Hence $\synval(\game) = \frac12$ while $1 - \valstar(\game) = 0$.

This game is not symmetric, and for this instance symmetry alone repairs the conclusion:
symmetrizing $D$ so as to accept every unequal answer pair off the diagonal gives
$\synval(\game) = 1$ too, by $P = \ket{0}\bra{0}$, $Q = \ket{1}\bra{1}$ in dimension $2$.
The counterexample therefore refutes the theorem as quantified over all synchronous games,
rather than exhibiting a failure among the symmetric games the pipeline produces; what the
source's proof needs in either case is the diagonal weight itself, the synchronicity
defect being measured against the marginal and controlled only by the diagonal pairs.

A positive diagonal weight does not rescue the printed statement either: the hypothesis
has to be quantitative. For $p \in (0,1)$ let $\mu$ give weight $p/2$ to each of $(1,1)$
and $(2,2)$ and $(1-p)/2$ to each of $(1,2)$ and $(2,1)$, with $D$ accepting equal answers
on the diagonal and $(a,b) = (0,1)$ off it; every question then carries diagonal weight
exactly $p$. Answering $0$ and $1$ deterministically loses only the diagonal, so
$\valstar(\game) \geq 1 - p$, while a synchronous strategy wins the diagonal with certainty
($\sum_a \tau(M^x_a M^x_a) = 1$ for a projective family) and, by the computation above, at
most half of the remaining weight, so $\synval(\game) = \frac{1+p}{2}$. Thus
$1 - \synval(\game) = \frac{1-p}{2}$ stays above $\frac14$ for $p \leq \frac12$ while
$1 - \valstar(\game) \leq p$ tends to $0$: no constants independent of the game satisfy the
printed inequality, and the $\kappa$ of the corrected one cannot be absorbed into $C$.

Nor does the hypothesis survive repetition: if $\mu(x,x) = \kappa \sum_y \mu(x,y)$ for
every $x$, the $k$-fold product distribution satisfies the same inequality with $\kappa^k$
and no better. A bound on the synchronous value of a $k$-fold repetition therefore says
nothing about its quantum value once $k$ grows with the instance size. For that reason
Theorem~\ref{thm:almost-sync} is used in this blueprint only on the game being repeated,
inside the proof of Theorem~\ref{thm:parallel-repetition}, never on a repeated game, and
the soundness of the pipeline is stated in $\valstar$ from that point onwards. The
invariant that supplies the hypothesis there is Remark~\ref{rem:sync-invariant}.

Even with the hypothesis, the transport cannot by itself reach the threshold
$\frac12$ that Definition~\ref{def:mipstar} demands. Take $\mX = \{1,2\}$, a singleton
answer alphabet, $\mu(1,1) = \mu(2,2) = \frac12$, and $D$ accepting at $(1,1)$ only: both
values are $\frac12$ and $\kappa = 1$, so the corrected inequality forces $C \geq
2^{c-1}$, whence $(2C)^{-1/c} \leq \frac12$. From $\synval(\game) \leq \frac12$ it then
concludes at best $\valstar(\game) \leq 1 - \kappa (2C)^{-1/c}$, and that bound is never
better than $\frac12$. The threshold must therefore come from
Theorem~\ref{thm:direct-repetition-q}, which is a $\valstar$ theorem, and carrying
soundness in $\valstar$ from the repetition step outwards is not a convenience but the
only available route.
\end{remark}

\subsection{The orthonormalization lemma}
\label{sec:rr-orthonormalization}

\effortMedium\ \qlib

\begin{theorem}[Orthonormalization]
\label{thm:orthonormalization}
\ledgernode{1.2.1.7.1, 1.2.1.7.1.1, 1.2.1.7.1.2, 1.2.1.7.1.3}
\uses{def:distance}
\lean{MIPRE.Orthonormalization.povm_orthogonalization_finDim}
\leanok
Let $\cal{M}$ be a von Neumann algebra on a finite-dimensional Hilbert space, let
$\varphi$ be a normal state on $\cal{M}$, and let $\{M_a\}_{a \in \mA}$ be a POVM in
$\cal{M}$ that is $\eps$-nearly projective:
\begin{equation*}
  \varphi\Bigl( \sum_{a \in \mA} M_a^2 \Bigr) > 1 - \eps\;.
\end{equation*}
Then there is a projective measurement $\{P_a\}_{a \in \mA}$ in $\cal{M}$ with
\begin{equation*}
  \varphi\Bigl( \sum_{a \in \mA} (M_a - P_a)^*(M_a - P_a) \Bigr) < 9\eps\;.
\end{equation*}
The same holds, on average, for a family $\{M^x_a\}$ indexed by $x \sim \mu$, by applying
it to $\varphi$ averaged over $x$.
\end{theorem}

\begin{remark}[What is proved, and in what generality]
\label{rem:orthonormalization-scope}
\uses{thm:orthonormalization}
\leanok
The statement above is the finite-dimensional case, which is the one the main theorem
needs and the one that is formalized. The vendored development proves three
unconditional cases, in increasing generality of the algebra and decreasing generality of
the state: every von Neumann algebra on a finite-dimensional space, for every normal
state (\texttt{MIPRE.Orthonormalization.povm\_orthogonalization\_finDim}, the
statement above, bridging \texttt{povm\_orthogonalization\_finDim\_vn}); $\cal{B}(H)$
for arbitrary $H$ (\texttt{povm\_orthogonalization\_fullAlgebra\_general}); and II$_1$
factors whose trace and state are both of trace-class form
(\texttt{povm\_orthogonalization\_II\textsubscript{1}Factor}), which is the case a
tracial, commuting-operator strategy lives in.

The general case --- every von Neumann algebra, every normal state --- is proved only as
an implication. Its hypothesis, \texttt{MvNStructureTheory}, bundles eight facts: seven
are classical Murray--von Neumann structure theory that Mathlib does not have (normal
functionals are of trace-class form; the type decomposition; the finite decomposition;
increasing nets of finite projections in the semifinite part; the center-valued trace;
halving in type III), and the eighth is the measurable-selection step of the source's own
Lemma~3.1 for the finite type~I part. No term of that type is constructed, so the
unconditional general statement remains open upstream, and with it the three companion
results (Theorems 1.1 and 1.4 and Corollary 1.5 of~\cite{dlS21}), which are formalized in
finite dimension and conditionally in general.
\texttt{MIPRE/Background/Orthonormalization/Axioms.lean} makes the split mechanical: it
fails to build if an unconditional result acquires an axiom, or if one of the four open
targets is closed upstream without this blueprint being updated.
\end{remark}

\paragraph{Source.} de la Salle~\cite{dlS21}, which gives a dimension-free bound with
explicit constants; earlier versions with worse parameters appear in~\cite{KPS18}
and in~\cite{JNVWY20}. The Lean proof is vendored from the
\texttt{orthogonalization} package of the \texttt{commuting-repetition} repository
(the same upstream as Theorem~\ref{thm:direct-repetition-co}) by
\texttt{scripts/vendor-repetition.py}; see that directory's \texttt{README.md} for what
the copy changes (import prefixes, nothing else).

\paragraph{Comments.} Used constantly: soundness analyses produce POVMs that are only
approximately projective, and this lemma replaces them by genuinely projective
measurements on the same space, which is what makes the ``projective strategies
throughout'' convention of Section~\ref{sec:departures} viable (avoiding Naimark
dilation). The finite-dimensional case suffices for the main theorem and is a clean,
self-contained formalization target in linear algebra.

The first of the two claims about the source made by the verification
campaign~\cite{Audit26} --- that this statement was over-hypothesized, the result needing
only a \emph{normal} state rather than a normal faithful tracial one, with constant $9$
and exponent $1$ --- is now settled in the direction the campaign predicted, and settled
by proof rather than by reading: the statement above carries neither faithfulness nor
traciality and is checked. That matters beyond bookkeeping, because the tracial
hypothesis excluded the bipartite vector-state applications that several of the
projectivization sites of Section~\ref{ch:proof-structure} need; those sites can now use
this lemma, where before they needed Naimark dilation or the projective-by-definition
convention of Section~\ref{sec:departures}. What is settled is the mathematics, not the
attribution: the campaign's second claim, that crediting the weaker predecessor
to~\cite{KPS18} and~\cite{JNVWY20} is wrong and the source credits the low individual
degree paper and Kempe--Vidick instead, is still unchecked here, \cite{dlS21} not being
vendored in this repository. The campaign also offers a consistency-to-projectivity step
giving $O(\delta)$ in place of $O(\delta^{1/4})$, which would sharpen the proof of
Theorem~\ref{thm:almost-sync}.

One difference from the statement this blueprint carried before: the hypothesis and the
conclusion are both strict, and the near-projectivity hypothesis is written
$\varphi(\sum_a M_a^2) > 1 - \eps$ rather than $\sum_a \tau(M_a - M_a^2) \leq \eps$. For a
POVM the two agree, since $\sum_a M_a = 1$; converting the strict form to a non-strict one
costs the constant, $9$ becoming $10$.

\subsection{Magic Square rigidity}
\label{sec:rr-magic-square}

\effortEasy\ \qlib

\begin{theorem}[Magic Square rigidity]
\label{thm:ms-rigidity}
\ledgernode{1.2.2.4.1}
\uses{def:q-value,def:distance}
Let $\strategy = (\ket{\psi}, A, B)$ be a strategy that succeeds in the Magic Square
game $\game^{\MS}$ with probability $1 - \eps$. Then there exist local isometries
$\phi_\alice : \mH_\alice \to (\C^2)^{\ot 2} \ot \mH_{\alice''}$, $\phi_\bob :
\mH_\bob \to (\C^2)^{\ot 2} \ot \mH_{\bob''}$ and a state $\ket{\aux} \in
\mH_{\alice''} \ot \mH_{\bob''}$ such that
\begin{equation*}
  \norm{ \phi_\alice \ot \phi_\bob \ket{\psi} - \ket{\EPR_2}^{\ot 2} \ot \ket{\aux} }
  \leq O(\sqrt{\eps})\;,
\end{equation*}
and, writing $\tilde A = \phi_\alice A \phi_\alice^\dagger$, the observables for the
questions $\Variable_1$ and $\Variable_5$ satisfy
$\tilde A^{\Variable_1}_b \approx_{\sqrt\eps} \sigma^X_b \ot \Id$ and
$\tilde A^{\Variable_5}_b \approx_{\sqrt\eps} \sigma^Z_b \ot \Id$ on
$\ket{\EPR_2}^{\ot 2} \ot \ket{\aux}$ (and symmetrically for $\tilde B$). In
particular the corresponding $\pm 1$-observables approximately anticommute:
$\tilde A^{\Variable_1} \tilde A^{\Variable_5} \approx_{\sqrt\eps}
- \tilde A^{\Variable_5} \tilde A^{\Variable_1}$.
\end{theorem}

\paragraph{Source.} Wu, Bancal, McKague and Scarani~\cite{WBMS16}; the game is the
Mermin--Peres magic square. The form above is~\cite[Section 8]{JNVWY20}.

\paragraph{Comments.} This is the elementary seed of all rigidity in the proof: it
certifies a single pair of anticommuting observables, which the low individual degree
test then amplifies to many qubits. The game is constant-sized, the proof is a page of
algebra, and it is an excellent first self-testing target. The synchronous restatement
(one measurement family, tracial state) should be derived once and reused as a template
for the other rigidity statements.

\subsection{The classical low individual degree test}
\label{sec:rr-lidt}

\effortHard

\begin{definition}[Low individual degree test]
\label{def:lidt}
\uses{def:game}
\lean{MIPRE.LIDT.lidtGame}
\leanok
Let $\F$ be a finite field with $q$ elements and let $m \geq 1$ and $d$ be integers.
The \emph{$(m,q,d)$-low individual degree test} is the two-player game in which the
verifier picks one of the following three sub-tests with probability $\tfrac13$ each,
having sampled a uniform point $u \sim \F^m$.
\begin{enumerate}
\item \textbf{Axis-parallel lines test.} Sample a uniform direction $i \sim [m]$ and a
  uniform role. One player receives the line $\ell = \{u + t e_i\}$ and answers a
  univariate polynomial $f$ of degree at most $d$; the other receives $u$ and
  answers $a \in \F$. Accept if $f(u) = a$.
\item \textbf{Self-consistency test.} Both players receive $u$ and answer elements of
  $\F$. Accept if the answers are equal.
\item \textbf{Diagonal lines test.} Sample a uniform $j \sim [m]$, a uniform direction
  $v \in \F^m$ whose coordinates after the $j$-th vanish, and a uniform role. One
  player receives the line $\ell = \{u + t v\}$ and answers a univariate polynomial
  of degree at most $md$; the other receives $u$ and answers $a \in \F$. Accept if
  $f(u) = a$.
\end{enumerate}
\end{definition}

\paragraph{Comments.} This is Figure~1 of~\cite{JNVWY21qld}. The direction $v$ of the
diagonal lines test may be zero, in which case the ``line'' is the single point $u$;
this degenerate case is part of the test as stated. In the formalization, line questions
are the lines themselves, presented canonically (the direction is rescaled so that its
first nonzero coordinate is $1$, and the base point is the point of the line whose
corresponding coordinate vanishes), so that the same line asked by two sub-tests is the
same question; a line answer is a polynomial in the parameter of that canonical
parametrization.

\begin{theorem}[Quantum soundness of the classical low individual degree test]
\label{thm:lidt-soundness}
\ledgernode{1.2.1}
\uses{def:lidt,def:tensor-strategy,def:inconsistency}
\lean{MIPRE.LIDT.lowIndividualDegree_soundness}
\leanok
Let
\begin{equation*}
  \delta_{lidt}(\eps, m, d, q, k)
  = 10^5 \, k^2 m^4 \bigl( \eps^{1/40000} + (d/q)^{1/40000}
      + e^{-k/(2560000 m^2)} \bigr).
\end{equation*}
Let $\strategy = (\ket{\psi}, A, B)$ be a tensor-product strategy for the
$(m,q,d)$-low individual degree test that succeeds with probability at least $1 - \eps$,
and let $k$ be an integer with $k \geq 400md$ and $k > 0$. Then there exist projective
measurements $G^\alice = \{G^\alice_g\}$ on Alice's space and $G^\bob = \{G^\bob_g\}$ on
Bob's space, both with outcomes the polynomials $g$ in $m$ variables of individual degree
at most $d$, such that, writing $\delta = \delta_{lidt}(\eps, m, d, q, k)$ and letting
$u \sim \F^m$ be uniform:
\begin{enumerate}
\item Alice's point measurement $A^u$ and the evaluation $G^\bob_{[g(u) = \cdot]}$
  are $\delta$-consistent;
\item the evaluation $G^\alice_{[g(u) = \cdot]}$ and Bob's point measurement $B^u$
  are $\delta$-consistent;
\item $G^\alice$ and $G^\bob$ are $\delta$-consistent.
\end{enumerate}
\end{theorem}
\begin{proof}
\leanok
This is Theorem~1.3 of~\cite{JNVWY21qld} (\texttt{thm:main-formal} there), whose proof
is the MIPStarRE formalization vendored into this repository. The Lean proof translates
a strategy for our game into MIPStarRE's strategy container (the state as a density
matrix; the line measurements as coarse-grainings of the measurements at the canonical
line questions along the change of parametrization, which makes them covariant under
rebasing), shows that MIPStarRE's failure surrogate is at most the rejection probability
$1 - \mathrm{val}$ (the consistency defect of each branch is at most the rejection
probability of that branch, since acceptance forces the coarse-grained answers to
agree), applies their main theorem, and reads the resulting measurements back into our
vocabulary.
\end{proof}

\paragraph{Source.} Ji, Natarajan, Vidick, Wright and Yuen~\cite{JNVWY21qld},
Theorem~1.3 (the formal version for general, not necessarily symmetric, strategies).
The proof is a formalization by Sirui Lu and collaborators
(\url{https://github.com/LionSR/MIPStarRE}), vendored into this repository with the
authors' permission at \texttt{MIPRE/Background/LIDT/MIPStarRE/}; the statement above is
this project's own, and is connected to theirs by an explicit bridge.

\paragraph{Comments.} Two hypotheses correct the printed statement, as established by
that formalization: the sampling parameter satisfies $k \geq 400md$ where the paper
prints $k \geq md$, and $k$ must be positive, since the printed error vanishes at
$k = 0$. The parameter $k$ remains free: the error carries both a factor $k^2$ and a
term $e^{-k/(2560000m^2)}$, so no single choice of $k$ reduces it to a function of
$(\eps, m, d, q)$; a corollary in the shape of $\delta_{qld}$ below is a separate
derivation.


\subsubsection{How Theorem~\ref{thm:lidt-soundness} is proved}

The proof is a reduction to a \emph{two-prover tensor code test}, and the reduction is the
paper's own; only the tensor code theorem is imported. This matters for reading the Lean
development, whose proof of Theorem~\ref{thm:lidt-soundness} follows the same shape
(Remark~\ref{rem:lidt-formalized}).

\effortMedium

\begin{lemma}[Setup, and the two tests are not the same test]
\label{lem:lidt-reduction-setup}
\ledgernode{1.2.1.1, 1.2.1.2}
\uses{def:lidt,def:tensor-strategy}
Let $C$ be the Reed--Solomon code of degree $d$ over $\F_q$. The codewords of $C^{\ot m}$ are
exactly the evaluation tables of $m$-variate polynomials of individual degree at most $d$. The
low individual degree test is \emph{not} the two-prover tensor code test --- the two differ in
which lines they sample and how they label answers --- so a strategy for the former must be
converted, and the converted strategy $\strategy'$ is well defined and projective: points by
reversing the coordinate order, lines by choosing a seed in the fibre of size $q/m$ with
answers relabelled through the parametrization of the line, pairs by restriction.
\end{lemma}

\effortHard

\begin{lemma}[Transfer of the line and pair tests]
\label{lem:lidt-test-transfer}
\ledgernode{1.2.1.3, 1.2.1.4}
\uses{lem:lidt-reduction-setup,def:lidt}
Assume $m \mid q$. Then a uniform axis-parallel line of $[q]^m$ together with a uniform seed
from its fibre and a uniform point on it induces exactly the axis-parallel-line distribution
the tensor code test uses; and for two independent uniform points of a subcube fixing the last
$j - 1$ coordinates, the pair (point, difference) is distributed as the diagonal-line test
requires, the seed law $\Pr[\chi(s) = i] \propto q^{i-1}$ matching the conditional law of the
low individual degree test.
\end{lemma}

\begin{proof}
The line case needs the conditional law of the seed given the line, which is immediate from
the joint distribution rather than part of the cited normal form; the divisibility $m \mid q$ is
where it is used, and is the same hypothesis that appears in
Remark~\ref{rem:pcp-power-of-two} for the same reason.
\end{proof}

\effortMedium

\begin{lemma}[Transfer of synchronicity]
\label{lem:lidt-sync-transfer}
\ledgernode{1.2.1.5}
\uses{lem:lidt-reduction-setup,def:inconsistency,lem:schwartz-zippel}
Points synchronicity of $\strategy'$ fails with probability at most $9\eps$, and lines
synchronicity with probability at most $81\eps + d/q$.
\end{lemma}

\begin{proof}
\uses{lem:schwartz-zippel}
A triangle chain through the other player's point measurement, plus the observation that two
distinct univariate polynomials of degree $d$ agree at no more than $d$ of the $q$ points,
which is Lemma~\ref{lem:schwartz-zippel} in one variable.
\end{proof}

\effortMedium

\begin{lemma}[Summation, derandomization, and the parameter choice]
\label{lem:lidt-derandomize}
\ledgernode{1.2.1.6, 1.2.1.8}
\uses{lem:lidt-test-transfer,lem:lidt-sync-transfer,thm:tensor-codes}
Summing the per-subtest bounds with the tensor code test's weights gives expected total failure
at most $\eps' = C(\eps + md/q)$ over the independent uniform choices of seeds, so some
deterministic choice achieves $\eps'$. For the Reed--Solomon code the number of codeword
coefficients is $t = d + 1$, and $k = m^3 d + 12m(d+1)$ satisfies the tensor code theorem's
hypothesis $k \geq 12mt$ for \emph{all} $m, d \geq 1$ while keeping $e^{-k/m^2} \leq e^{-md}$ and
$k = \poly(m, d)$.
\end{lemma}

\effortHard

\begin{theorem}[Two-prover tensor code test]
\label{thm:tensor-codes}
\ledgernode{1.2.1.7, 1.2.1.7.3}
\uses{def:tensor-strategy,def:inconsistency,thm:orthonormalization,thm:almost-sync}
A two-prover strategy passing the two-prover tensor code test for $C^{\ot m}$ with probability
$1 - \eps'$, with the free parameter $k \geq 12mt$, is close to one measuring codewords of
$C^{\ot m}$, with error
$\eta = \poly(m, t, k) \cdot \poly\bigl( \eps', n^{-1}, e^{-\Omega(k/m^2)} \bigr)$.
\end{theorem}

\begin{proof}
Imported from~\cite{JNVWY21tensor}, in the published form. The source is vendored and was
audited in full; the statement above is its synchronous form instantiated at the free
parameter, and Lemma~\ref{lem:tensor-codes-bipartite} is the two-prover form the reduction
actually consumes.
\end{proof}

\effortHard

\begin{lemma}[The induction inside the tensor code theorem]
\label{lem:tensor-codes-induction}
\ledgernode{1.2.1.7.4, 1.2.1.7.5}
\uses{thm:tensor-codes,lem:tensor-codes-tools}
The theorem is proved by induction on the number of tensor factors. For a single factor a
synchronous strategy passing the test yields a measurement over codewords of $C$ with the
claimed consistency. The step from $m$ to $m+1$ composes the two tools of
Lemma~\ref{lem:tensor-codes-tools}, and the closing arithmetic propagates the error terms with
universal constants.
\end{lemma}

\effortHard

\begin{lemma}[Self-improvement and pasting]
\label{lem:tensor-codes-tools}
\ledgernode{1.2.1.7.6, 1.2.1.7.7}
\uses{thm:orthonormalization,def:distance}
Two tools drive the induction. \emph{Self-improvement}: a measurement family that almost
satisfies the test's conditions can be replaced by an exactly projective family that is
approximately consistent, at controlled error. \emph{Pasting}: measurements for the first $j$
coordinates and for coordinate $j+1$ that approximately commute on average --- which the
subcube commutation test certifies --- can be pasted into a joint measurement over the
restricted tensor code.
\end{lemma}

\effortHard

\begin{lemma}[From the synchronous to the two-prover form]
\label{lem:tensor-codes-bipartite}
\ledgernode{1.2.1.7.8}
\uses{thm:tensor-codes,thm:almost-sync}
The two-prover form of Theorem~\ref{thm:tensor-codes} follows from the synchronous form by
symmetrizing the strategy --- the swap construction, which preserves the passing probability
--- and then applying Theorem~\ref{thm:almost-sync} with a concavified error function.
\end{lemma}

\begin{remark}[Almost-synchronicity is off the pipeline but inside an import]
\label{rem:almost-sync-transitive}
\uses{thm:almost-sync,lem:tensor-codes-bipartite,thm:lidt-soundness,rem:almost-sync-hypothesis}
This is worth stating plainly, because the blueprint elsewhere says
Theorem~\ref{thm:almost-sync} is not needed, and that is only true of the pipeline. The
pipeline does not invoke it: soundness is carried in $\valstar$ from end to end, so no
transport from $\synval$ is required (Remark~\ref{rem:almost-sync-hypothesis} explains why its
hypothesis would not even be available at the outer end). But
Lemma~\ref{lem:tensor-codes-bipartite} consumes it, and
Theorem~\ref{thm:lidt-soundness} reduces to that lemma --- so the blueprint depends on
almost-synchronicity \emph{transitively, through an import}, and cannot simply drop it.

What makes this benign rather than alarming is that the dependency is already discharged in
Lean. The proof of \texttt{MIPRE.LIDT.lowIndividualDegree\_soundness} is sorry-free
(Remark~\ref{rem:lidt-formalized}), so everything its chain needs --- this included --- is
proved from Mathlib in the vendored development. The practical consequence is for planning
rather than for soundness: issues on the synchronous transport are about the
\emph{pipeline's} use of Theorem~\ref{thm:almost-sync}, which is genuinely optional, and not
about this one, which is not.
\end{remark}

\effortHard

\begin{lemma}[Several codewords: the simultaneity reduction]
\label{lem:lidt-ldc}
\ledgernode{1.2.1.9, 1.2.1.9.1, 1.2.1.9.2, 1.2.1.9.3}
\uses{lem:lidt-derandomize,def:lidt}
The general case of Theorem~\ref{thm:lidt-soundness}, with several codewords tested at once,
reduces to the single-codeword case \emph{in the paper}, importing nothing further. The
reduction replaces the parameters $(q, m, d, \mathrm{ldc})$ by $(q, M, d, 1)$ where $M$ is the
least power of two with $M \geq m + \mathrm{ldc}$ --- forced by the divisibility $M \mid q$, and
satisfying $M \leq 4m\,\mathrm{ldc}$ --- leaving the degree unchanged. The resulting error is
$\delta_{lidt}(\eps, q, m, d, \mathrm{ldc}) = a (dm\,\mathrm{ldc})^a \bigl( \eps^b + q^{-b} +
2^{-bmd} \bigr)$, the $\mathrm{ldc}^a$ factor and the substitution $m \mapsto M$ being the only
changes.
\end{lemma}

\begin{proof}
The construction is modelled on the corresponding derivation in~\cite{NW19}, but is written out
for this test rather than cited: the vendored audit of that source confirmed that the theorem
it would have imported is not the one needed here.
\end{proof}

\subsection{The quantum low individual degree test}
\label{sec:rr-qld}

\effortHard\ \qlib

\begin{remark}[The one background result that is already formalized]
\label{rem:lidt-formalized}
\ledgernode{1.2.1, 1.2.1.7.2.4}
\uses{thm:lidt-soundness,def:lidt}
Theorem~\ref{thm:lidt-soundness} is the only statement of this section whose Lean proof is
complete, and the way it got there is worth recording, because the ledger and the
formalization disagree about what kind of statement it is.

The paper proves node $1.2.1$ by an explicit reduction --- nodes $1.2.1.1$ to $1.2.1.8$ ---
to the two-prover tensor-code test of~\cite{JNVWY21tensor}, and the general case of several
codewords is proved in the paper too, importing nothing. So the assumption is narrower than
the statement: exactly one node, $1.2.1.7$, is an import, namely the tensor-code theorem
itself. The Lean statement is not an axiom, and
its \texttt{\#print axioms} guard fails the build if it ever becomes one: it is proved
from Mathlib alone through the vendored formalization of~\cite{JNVWY21qld}, with the
bridge modules of \texttt{MIPRE/Background/LIDT/Bridge/} translating this development's
games, strategies and consistency relation into that one's and back. So on this node the
formalization is strictly stronger than the paper, and the $28$ nodes of the subtree are
accounted for as the internal structure of a proof that is finished.

It also corrected the statement. The printed parameter constraint is $md \leq k$; the
formalization establishes it only for $400\,md \leq k$, and additionally needs $k > 0$.
Both corrections are in the statement above, and they cost nothing downstream because $k$
is free --- the error carries a factor $k^2$ against a term $e^{-k/(2560000 m^2)}$, so it
is minimized at some finite $k$ either way.

One node of the subtree is worth reading before reusing anything from it. Node
$1.2.1.7.2.4$ records that the interface between the product-basis transposes
of~\cite{JNVWY20} and the Schmidt-diagonal symmetric strategies of~\cite{Vid22} is
\emph{false} as stated, with an explicit $2 \times 2$ counterexample. It is closed, not
open: the repair is an approximate Takagi bridge, which puts every strategy of the first
form within a controlled synchronicity defect of one of the second after a local basis
change, with the closure constants re-derived and checked by a registered claim-test. It
does not affect Theorem~\ref{thm:lidt-soundness}, whose Lean proof does not go through the
interface at all. The lesson for a formalization is that the bridge is the thing to state,
not the exact identity it replaces.
\end{remark}

\begin{theorem}[Pauli basis test / quantum low individual degree test]
\label{thm:qld}
\ledgernode{1.2.2, 1.2.2.17}
\uses{def:q-value,def:distance,thm:ms-rigidity,thm:lidt-soundness}
There exists a function
\begin{equation*}
  \delta_{qld}(\eps, m, d, q) = a (md)^a \bigl( \eps^b + q^{-b} + 2^{-bmd} \bigr)
\end{equation*}
for universal constants $a \geq 1$ and $0 < b < 1$ such that the following holds. For
all admissible parameter tuples $\qldparams = (q, m, d)$ and all strategies
$\strategy = (\ket{\psi}, A, B)$ for the game $\game^{\pauli}_{\qldparams}$ that
succeed with probability at least $1 - \eps$, there exist local isometries
$\phi_\alice, \phi_\bob$ mapping into $\mH' \ot (\C^q)^{\ot M}$ with $M = 2^m$, and a
state $\ket{\aux}$, such that
$\norm{\phi_\alice \ot \phi_\bob \ket{\psi} - \ket{\aux} \ot \ket{\EPR_q}^{\ot M}}
\leq \delta_{qld}$, and for $W \in \{X, Z\}$ the measurements for the questions
$(\pauli, W)$ are $\delta_{qld}$-close to the Pauli basis measurements $\tau^W_u$, $u
\in \F_q^M$, acting on $\ket{\EPR_q}^{\ot M}$.
\end{theorem}

\paragraph{Source.} Ji, Natarajan, Vidick, Wright and Yuen~\cite{JNVWY21qld}, which
proves quantum soundness of the classical low individual degree test and derives the
Pauli basis test; this corrects the analysis of the total degree test relied on in
early versions of~\cite{JNVWY20}. The reduction from the Pauli test to the low
individual degree test follows~\cite[Section 7]{JNVWY20}, with the Magic Square
theorem (Theorem~\ref{thm:ms-rigidity}) providing the base anticommutation.

\paragraph{Comments.} This is the largest and hardest imported result: its proof is an
induction on the number of variables with a delicate self-improvement argument. It is
the workhorse of question reduction (introspection certifies the players' sampling of
their own questions against Pauli basis measurements). Its classical ingredient is now
available as Theorem~\ref{thm:lidt-soundness}; what remains is the reduction from the
Pauli basis test, which uses Theorem~\ref{thm:ms-rigidity} for the base anticommutation
and Theorem~\ref{thm:gowers-hatami} and Theorem~\ref{thm:orthonormalization} as black
boxes, together with a corollary of Theorem~\ref{thm:lidt-soundness} collapsing the
sampling parameter $k$ into a single error function.

The statement is bipartite, so every use of it inside the synchronous pipeline costs one
application of the transport of Theorem~\ref{thm:almost-sync}, with its hypothesis and its
loss $\eps \mapsto \eps^{c}$; restating the rigidity lemmas tracially once, rather than
transporting at each use, is the cheaper route and is not yet done. The verification
campaign~\cite{Audit26} further reports three corrections in the tensor-code material the
proof passes through --- $\kappa$-convexity requires $c_1 \leq 1$, the synchronicity
defect obeys $\delta_{\mathrm{sync}} \leq 4\eps$ rather than $3\eps$, and one corner block
pairs with a transpose --- none checked here, none affecting the statement above, all
relevant to whoever writes its proof.


\subsubsection{The shape of the soundness proof}

\begin{remark}[How the soundness proof of Theorem~\ref{thm:qld} is organized]
\label{rem:qld-architecture}
\uses{thm:qld,thm:lidt-soundness,thm:orthonormalization,thm:ms-rigidity}
The soundness half of Theorem~\ref{thm:qld} is not proved in the main text of~\cite{JNVWY21qld};
it is a five-file appendix, and it is the largest single proof the pipeline imports. Its
architecture is worth stating before its pieces, because the pieces are unreadable without it.
There are five stages.
\begin{enumerate}
\item \textbf{Preliminaries}: three utility lemmas about closeness, one of which
  (Corollary~\ref{cor:ortho-from-consistency}) is used on almost every page.
\item \textbf{Win implications}: unpack ``succeeds with probability $1 - \eps$'' into one
  consistency or commutation statement per subtest (Lemma~\ref{lem:qld-win}).
\item \textbf{Expansion}: adjoin EPR ancillas and replace the strategy's measurements by ones
  on the enlarged space that exactly commute across the two bases
  (Lemmas~\ref{lem:qld-expanded-points} and~\ref{lem:qld-expanded-lines}).
\item \textbf{Combining}: build a \emph{single} measurement returning both the $X$- and the
  $Z$-polynomial, by padding the ambient space and applying the classical low individual
  degree test to the padded points (Lemma~\ref{lem:qld-simultaneous}).
\item \textbf{Separation and the swap isometry}: extract exact Pauli operators from that
  single measurement and rotate them onto honest EPR pairs
  (Lemmas~\ref{lem:qld-exact-paulis} and~\ref{lem:qld-swap}).
\end{enumerate}
Two features of this decomposition matter for a formalization. Stage 4 \emph{consumes}
Theorem~\ref{thm:lidt-soundness}: the simultaneous measurement is obtained by applying the
classical test's quantum soundness to the padded point measurements with parameters
$(q, 4m, d)$ and a single codeword. So the one background result whose Lean proof is already
finished and sorry-free (Remark~\ref{rem:lidt-formalized}) is an input to the one that is the
largest remaining gap --- which is the right way round, and means the hard analytic core of
stage 4 is done. And stage 5, not stage 4, is where the exactness comes from: the Pauli group
relations are recovered \emph{exactly}, not approximately, which is what lets the swap
isometry be the standard one.
\end{remark}

\subsubsection{Preliminaries}

\effortEasy

\begin{lemma}[Averaging preserves closeness]
\label{lem:qld-averaging}
\ledgernode{1.2.2.1}
\uses{def:distance}
Let $A^x \approx_\delta B^x$ on average over $x \sim \mu$, and let $\alpha_x$ be coefficients
with $\abs{\alpha_x} \leq 1$. Then $\E_x \alpha_x A^x \approx_\delta \E_x \alpha_x B^x$.
\end{lemma}

\begin{proof}
Triangle inequality and Jensen. Note the hypothesis is $\abs{\alpha_x} \leq 1$, not
$\abs{\alpha_x} = 1$: the lemma is used with genuinely contractive coefficients.
\end{proof}

\effortEasy

\begin{lemma}[From POVM elements to generalized observables]
\label{lem:qld-povm-to-obs}
\ledgernode{1.2.2.2}
\uses{def:distance,lem:qld-averaging}
Let $\{A^x_a\}$ and $\{B^x_a\}$ be POVMs with a common outcome set $\mA$, and let $\alpha_a$
be coefficients of unit modulus. If $A^x_a \approx_\delta B^x_a$ then the generalized
observables $A^x = \sum_a \alpha_a A^x_a$ and $B^x = \sum_a \alpha_a B^x_a$ satisfy
$A^x \approx_{\abs{\mA} \delta} B^x$.
\end{lemma}

\effortMedium

\begin{corollary}[Orthonormalization from consistency]
\label{cor:ortho-from-consistency}
\ledgernode{1.2.2.3, 1.2.2.3.1}
\uses{thm:orthonormalization,def:inconsistency,def:distance}
There is $\eta(\delta) = O(\delta)$ such that the following holds. Let $\{Q_a\}$ be a POVM on
$\mH_\alice$ and $\{R_a\}$ a POVM on $\mH_\bob$ that are $\delta$-consistent on
$\ket{\psi}$. Then there is a \emph{projective} measurement $\{P_a\}$ on $\mH_\alice$ with
$P_a \ot \Id \approx_{\eta(\delta)} Q_a \ot \Id$ on $\ket{\psi}$.
\end{corollary}

\begin{proof}
\uses{thm:orthonormalization}
This is not Theorem~\ref{thm:orthonormalization} restated: that theorem's hypothesis is
near-projectivity, $\sum_a \tau(Q_a - Q_a^2) \leq \eps$, while what the appendix has is
bipartite consistency. The bridge is a Cauchy--Schwarz step. Write
$\gamma = \sum_{a \neq b} \bra{\psi} Q_a \ot R_b \ket{\psi}$, so that
$\sum_a \bra{\psi} Q_a \ot R_a \ket{\psi} = 1 - \gamma$. Since
$\sum_a R_a^2 \leq \sum_a R_a = \Id$, Cauchy--Schwarz over $a$ gives
$\sum_a \bra{\psi} Q_a^2 \ot \Id \ket{\psi} \geq (1 - \gamma)^2 \geq 1 - 2\gamma$, which is
near-projectivity with $\eps = 2\gamma$. Theorem~\ref{thm:orthonormalization}, applied to the
von Neumann algebra $\cal{L}(\mH_\alice) \ot \Id$ with the normal vector state of
$\ket{\psi}$, then returns the projective family, necessarily of the form $P_a \ot \Id$.
\end{proof}

\subsubsection{What winning implies}

\effortHard

\begin{lemma}[Win implications]
\label{lem:qld-win}
\ledgernode{1.2.2.4, 1.2.2.4.2}
\uses{thm:qld,thm:ms-rigidity,def:inconsistency,def:distance}
Assume $6md \leq q$. A strategy succeeding with probability $1 - \eps$ in the Pauli basis test
satisfies, with the conditional weight of each subtest carried explicitly, a list of
consistency and commutation relations: self-consistency of the point measurements; the
low-degree subtest's relations; the anticommutation relations coming from the Magic Square
subtest; the relations tying the $(\mathtt{Point}, W)$ measurements to the
$(\mathtt{Line}, W)$ ones; and Pauli-basis consistency, commutation, and
commutation-consistency.
\end{lemma}

\begin{proof}
\uses{thm:ms-rigidity}
Each item is the definition of the value of the corresponding subtest, divided by the
probability that the subtest is selected. The only non-elementary input is the
anticommutation item, which is where Theorem~\ref{thm:ms-rigidity} --- the admitted node --- is
consumed, and it is consumed exactly once.
\end{proof}

\subsubsection{Expansion: making the two bases commute exactly}

\effortHard

\begin{lemma}[Expanded point measurements]
\label{lem:qld-expanded-points}
\ledgernode{1.2.2.5}
\uses{lem:qld-win,cor:ortho-from-consistency,def:distance}
Adjoin EPR ancilla registers to $\ket{\psi}$. There are measurements
$\widehat{M}^{(\mathtt{Point}, W), u}$ on each player's space together with its ancilla such
that, with $\delta_{\mathrm{acom}}(\eps) = O(\sqrt{\eps})$: they are self-consistent across
the two sides at error $\eps$, and the $X$-side and $Z$-side coarse observables
\emph{approximately commute} on the expanded state at error $\delta_{\mathrm{acom}}$.
\end{lemma}

\begin{proof}
\uses{lem:qld-win}
The pivotal step is an exact cancellation of a sign $(-1)^\gamma$ between the strategy's own
observables, supplied by the anticommutation relations of Lemma~\ref{lem:qld-win}, and the
ancilla observables. Without the ancillas the two bases only commute up to an error that does
not improve; with them the sign cancels identically and what is left is the stated
$\sqrt{\eps}$.
\end{proof}

\effortHard

\begin{lemma}[Expanded line measurements]
\label{lem:qld-expanded-lines}
\ledgernode{1.2.2.6}
\uses{lem:qld-expanded-points,cor:ortho-from-consistency}
For each $W \in \{X, Z\}$ and each line $\ell$ there is a \emph{projective} measurement
$\widehat{M}^{(\mathtt{Line}, W), \ell}_f$, with outcomes the degree-$md$ polynomials on
$\ell$, such that with $\delta_{\mathrm{Line}}(\eps) = O(\sqrt{\eps})$ it is self-consistent
over the line--point distribution and consistent with the expanded point measurements of
Lemma~\ref{lem:qld-expanded-points}, in both the form that compares a line outcome with a
point outcome and the form that compares the evaluation of the line polynomial at the point.
\end{lemma}

\begin{proof}
Projectivity comes from an orthogonality property of the convolution defining the family,
not from Corollary~\ref{cor:ortho-from-consistency}; consistency in both forms is
Lemma~\ref{lem:qld-expanded-points} together with the line subtest's relations.
\end{proof}

\subsubsection{Combining the two bases}

The point of this stage is a single measurement that returns the $X$-polynomial and the
$Z$-polynomial at once. It is reached by padding the ambient space from $\F_q^m$ to
$\F_q^{4m}$, so that a point of the padded space encodes a point in each basis together with
a pair of coefficients, and then applying the classical low individual degree test to the
padded points.

\effortHard

\begin{theorem}[Exact linearity from approximate linearity]
\label{thm:linearity}
\ledgernode{1.2.2.7, 1.2.2.7.1, 1.2.2.7.2, 1.2.2.7.3}
\uses{def:distance,thm:gowers-hatami}
Let $\{O^u\}_{u \in \F_2^t}$ be observables satisfying approximate linearity on average,
$O^u O^{u'} \approx_\delta O^{u + u'}$. Then on an extended space there are \emph{exactly}
linear observables $\{L^u\}$ --- a representation of $\F_2^t$, with $L^0 = \Id$ --- such that
$L^u \approx_\delta O^u$ on average over $u$. The error is linear in $\delta$, with no loss in
the exponent.
\end{theorem}

\begin{proof}
Imported from Natarajan and Vidick~\cite{NV17}, whose proof is the Fourier argument: the
operators $\widehat{A}^u = \E_a (-1)^{a \cdot u} A(a)$ have squares forming a POVM by
Parseval, Naimark dilates that POVM to a projective family on an extended space, and the
exactly linear observables are recovered from it by inverting the transform. The source is
vendored and was audited line by line, including its state-dependent distance toolkit; the
only facts it needs beyond that are Parseval and Naimark.
\end{proof}

\effortHard

\begin{lemma}[Combined $XZ$ point measurements]
\label{lem:qld-combined-points}
\ledgernode{1.2.2.8}
\uses{thm:linearity,lem:qld-expanded-points,def:distance}
There is $\delta_Q(\eps) = \poly(\eps)$ and, for every pair $(x, z) \in \F_q^m \times \F_q^m$,
a projective measurement $\widehat{Q}^{x, z}_{a, b}$ jointly returning the $X$-value at $x$
and the $Z$-value at $z$, self-consistent on average over uniform $(x, z)$ and consistent
with \emph{both} ordered products $\widehat{M}^Z_b \widehat{M}^X_a$ and
$\widehat{M}^X_a \widehat{M}^Z_b$ of the expanded point measurements. It is built as a Fourier
transform of the exactly linear family supplied by Theorem~\ref{thm:linearity}. Its marginals
are not claimed here; they are obtained in Lemma~\ref{lem:qld-pairs-of-lines}.
\end{lemma}

\effortHard

\begin{lemma}[Pairs of lines]
\label{lem:qld-pairs-of-lines}
\ledgernode{1.2.2.9}
\uses{lem:qld-combined-points,lem:qld-expanded-lines}
There is $\delta_P(\eps, m, d, q) = \poly(\eps, md/q)$ and, for every pair of lines
$(\ell_X, \ell_Z)$, a POVM with outcomes pairs of degree-$md$ polynomials, consistent on
average over the pair-of-lines distribution with the combined point measurements evaluated at
points of the respective lines.
\end{lemma}

\effortMedium

\begin{lemma}[Padded point measurements]
\label{lem:qld-padded-points}
\ledgernode{1.2.2.10}
\uses{lem:qld-combined-points}
Write a point of the padded space $\F_q^{4m}$ as $u = (x, z, \alpha, \beta, w)$. The
measurement $\widehat{Q}^u$ returning the combined value $\alpha g_X(x) + \beta g_Z(z)$ is
well defined --- independent of the dummy coordinates $w$ --- projective, self-consistent, and
consistent with both ordered products of the expanded point measurements. The coarse-graining
onto the fibres $\{\alpha a + \beta b = c\}$ is justified by character orthogonality.
\end{lemma}

\effortMedium

\begin{lemma}[The sublines distribution]
\label{lem:qld-sublines}
\ledgernode{1.2.2.11}
\uses{def:lidt}
There is a distribution over triples $(\ell, \ell_X, \ell_Z)$ --- a line in the padded space
together with a line in each unpadded copy --- whose $\ell$-marginal is the line marginal of
the padded space's line--point distribution, whose $(\ell_X, \ell_Z)$-marginals are mixtures
of the restricted line--point distributions, and for which the evaluation constraints tying
$f|_\ell$ to the pair $(f_X, f_Z)$ hold pointwise.
\end{lemma}

\begin{proof}
Purely combinatorial: the three constraints are checked separately, and nothing quantum
enters. This is the most self-contained statement in the appendix and a reasonable place to
start formalizing it.
\end{proof}

\effortHard

\begin{lemma}[Padded line POVM]
\label{lem:qld-padded-lines}
\ledgernode{1.2.2.12}
\uses{lem:qld-pairs-of-lines,lem:qld-sublines,lem:qld-padded-points}
There is $\delta_{\mathrm{combine}}(\eps, m, d, q) = m \cdot \poly(\eps, md/q)$ and, for every
line $\ell$ in the padded space, a POVM with outcomes the degree-$(md + 1)$ polynomials on
$\ell$, consistent on average over the padded line--point distribution with the padded point
measurement evaluated at a point of $\ell$.
\end{lemma}

\effortHard

\begin{lemma}[The simultaneous global measurement]
\label{lem:qld-simultaneous}
\ledgernode{1.2.2.13}
\uses{lem:qld-padded-lines,lem:qld-padded-points,thm:lidt-soundness}
There is
$\delta_S(\eps, m, d, q) = a (md)^a \bigl( \eps^b + q^{-b} + 2^{-bmd} \bigr)$ and a projective
measurement $\widehat{S}_{g_X, g_Z}$ with outcomes pairs of polynomials of individual degree
at most $d$ on $\F_q^m$, such that for each $W \in \{X, Z\}$ the $W$-marginal evaluated at a
point is consistent with the expanded $(\mathtt{Point}, W)$ measurement, in the sandwiched
form $\widehat{M}^X \widehat{M}^Z \widehat{M}^X$ and its mirror.
\end{lemma}

\begin{proof}
\uses{thm:lidt-soundness}
Apply Theorem~\ref{thm:lidt-soundness} --- the quantum soundness of the \emph{classical} test
--- to the padded point and line measurements, with parameters $(q, 4m, d)$ and a single
codeword. This is the step at which the finished Lean proof of that theorem becomes an input
to this one (Remark~\ref{rem:qld-architecture}), and it is why $\delta_S$ has the shape of
$\delta_{lidt}$ rather than of the errors above it.
\end{proof}

\subsubsection{Separation and the swap isometry}

\effortMedium

\begin{lemma}[The marginals track the point measurements]
\label{lem:qld-helper}
\ledgernode{1.2.2.14}
\uses{lem:qld-simultaneous,def:distance}
For $W \in \{X, Z\}$ and on average over $u \in \F_q^m$, the $W$-marginal
$\widehat{S}^W_g$ of the simultaneous measurement satisfies
$\sum_g \widehat{S}^W_g \ot \widehat{M}^{(\mathtt{Point}, W), u}_{g(u)} \approx_{\delta_S} \Id$
and $\widehat{S}^W_g \bigl( \Id - \widehat{M}^{(\mathtt{Point}, W), u}_{g(u)} \bigr)
\approx_{\delta_S} 0$.
\end{lemma}

\effortHard

\begin{lemma}[Exact Pauli operators]
\label{lem:qld-exact-paulis}
\ledgernode{1.2.2.15}
\uses{lem:qld-simultaneous,lem:qld-helper,lem:qld-win,lem:schwartz-zippel}
From the simultaneous measurement one constructs, for $W \in \{X, Z\}$, measurements and
binary observables on the expanded space which satisfy the Pauli group relations
\emph{exactly} --- linearity in the argument, and $XZ$ commutation or anticommutation with the
correct phase --- and which are $\delta$-close to the strategy's
$(\mathtt{Point}, W)$ measurements on the expanded state.
\end{lemma}

\begin{proof}
\uses{lem:schwartz-zippel}
Exactness of the linearity comes from projectivity of the simultaneous measurement together
with the group law of the Pauli basis; Theorem~\ref{thm:linearity} is \emph{not} used here.
The remaining ingredient bounds the mass the simultaneous measurement places on
non-multilinear outcomes, via Lemma~\ref{lem:schwartz-zippel} applied in the correct
direction --- distinct polynomials of degree $md$ \emph{agree} on at most an $md/q$ fraction,
so it is the agreement and not the disagreement that the bound controls
(Remark~\ref{rem:qld-admitted}).
\end{proof}

\effortHard

\begin{lemma}[The swap isometry]
\label{lem:qld-swap}
\ledgernode{1.2.2.16}
\uses{lem:qld-exact-paulis,def:distance}
There is $\delta_{qld} = O\bigl( \delta_S^{1/4} + md/q \bigr)$ and local unitaries $V_\alice$,
$V_\bob$ on the expanded spaces such that $V_\alice \ot V_\bob$ maps the expanded state within
$\delta_{qld}$ of $\ket{\aux} \ot \ket{\EPR_q}^{\ot M}$, and conjugation by them carries the
exact Pauli observables of Lemma~\ref{lem:qld-exact-paulis} to the honest Pauli operators
acting on the EPR registers.
\end{lemma}

\begin{proof}
The standard swap-isometry argument. It is standard \emph{because} the representation is
exact: an approximate representation would not give a unitary, which is the whole reason
Lemma~\ref{lem:qld-exact-paulis} is stated with exact relations.
\end{proof}

\subsection{Undecidability of the halting problem}
\label{sec:rr-halting}

\effortMathlib

\begin{theorem}[Turing]
\label{thm:halting-undecidable}
\uses{def:re}
The halting problem is r.e.\ but not decidable; it is $\RE$-complete under many-one
(indeed polynomial-time) reductions.
\end{theorem}

\paragraph{Source.} Turing (1936); in Mathlib as part of the
\texttt{Computability} library.

\paragraph{Comments.} Listed to close the dependency graph: combined with the halting
reduction (Theorem~\ref{thm:halting}) it yields undecidability of approximating the
game value, and $\RE$-completeness gives $\RE \subseteq \MIPstar$.

\subsection{Cook--Levin and succinct satisfiability}
\label{sec:rr-cook-levin}

\effortHard

\begin{remark}[What stage 1.2 leaves open, and where its weight is]
\label{rem:qld-admitted}
\uses{thm:qld,thm:ms-rigidity,thm:lidt-soundness,def:distance}
Stage $1.2$ of the ledger is by far its heaviest: $62$ nodes and $169$ of the campaign's
$291$ challenges, against $46$ for answer reduction and $19$ for repetition. Three facts
about it bear on how this section should be formalized.

First, Theorem~\ref{thm:ms-rigidity} is one of only \emph{three} admitted nodes in the
whole campaign (node $1.2.2.4.1$; the others are anchored repetition and self-dual normal
bases). It is an import of~\cite[Theorem 6.9]{CS17}, not vendored and not reproved, with
two adaptations made in the paper: a weakening of trace-norm $O(\eps)$ closeness to
Euclidean $O(\sqrt{\eps})$, and a local basis change between two anticommuting pairs. The
saving grace for a formalization is how narrowly it is consumed --- only through one
anticommutation consequence, and not at all by the completeness leg, which uses the
in-file argument that anticommuting observables give a perfect Magic Square strategy
(Lemma~\ref{lem:mermin-peres} below is that strategy, and it is already in Lean).

Second, the soundness proof of Theorem~\ref{thm:qld} is not in the main text at all: it is
deferred to a five-file appendix, which is where most of stage $1.2$'s challenge weight
sits, including one of the campaign's three \emph{critical} findings, at node $1.2.2.15$:
the construction of exact Pauli operators applied Schwartz--Zippel in the wrong direction
--- distinct polynomials of degree $md$ \emph{agree} on at most an $md/q$ fraction, so for a
non-multilinear $g$ the disagreement indicator has expectation near $1$ --- which made the
step equivalent to an unproven assertion about how much mass the simultaneous measurement
places on non-multilinear outcomes. It is resolved, by bounding that mass first. A
formalization should treat the appendix as the real content and the main-text statement as
its interface.

The other two critical findings are worth naming here, because neither reaches this
blueprint and the reason is the same in both cases. At node $1.4$, the anchored repetition
theorem was applied to a game that is not of the anchored product form at all: detyping
destroys the product structure, and restricting a repeated strategy to the event that every
coordinate is a real question costs $16^k$. At node $1.5$, the compression statement had
omitted the doubly exponential term in its entanglement lower bound, which the recursion
needs. Both are resolved upstream; both live in the entanglement-form apparatus that
Theorem~\ref{thm:direct-repetition-q} and the $\valstar$ chain replace
(Remark~\ref{rem:admitted-nodes}). All three criticals are resolved, so this is not a
soundness worry --- it is an indication of where the delicate machinery is, and two thirds
of it is machinery this route does not carry.

Third, the introspection statement that consumes all of this (node $1.2$,
Theorem~\ref{thm:introspection}) leaves a soundness gap of $1/\poly(n)$ --- \emph{not} a
constant, which is the stronger reading and the one the ledger is explicit about. That is
exactly why parallel repetition has to follow it: repetition is what turns a gap shrinking
with $n$ back into the constant $\frac12$ the recursion needs.
\end{remark}

\begin{remark}[The three admitted nodes, and what this route does to them]
\label{rem:admitted-nodes}
\uses{thm:ms-rigidity,lem:self-dual-basis,thm:direct-repetition-q,rem:direct-vs-anchored}
With the ledger fully accounted for, the assumptions are countable. Of its $123$ nodes,
$120$ are validated and exactly three are \emph{admitted} --- asserted on the strength of a
cited source rather than proved. All three are now named in this blueprint, and the
value-form route treats them very differently.

\begin{enumerate}
\item \textbf{Anchored parallel repetition} (node $1.4.2$, the theorem
  of~\cite{BVY17}, which this blueprint no longer states). Not used.
  Direct repetition (Theorem~\ref{thm:direct-repetition-q}) replaces it, is formalized end
  to end, and takes a value hypothesis rather than an entanglement one
  (Remark~\ref{rem:direct-vs-anchored}). So one of the three admitted nodes leaves the
  pipeline altogether --- which is the strongest form of the argument that the
  value-form architecture is cheaper, not merely different.
\item \textbf{Magic Square rigidity} (node $1.2.2.4.1$, Theorem~\ref{thm:ms-rigidity}).
  Used, but narrowly: one anticommutation consequence, and not at all by the completeness
  leg, whose in-file argument is already in Lean (Lemma~\ref{lem:mermin-peres}).
\item \textbf{Efficient self-dual normal bases} (node $1.1.6.1$,
  Lemma~\ref{lem:self-dual-basis}). Unavoidable, and the only one of the three that is
  genuinely load-bearing here. It is also the most benign: three classical
  computational-algebra results, and the reason the admissible field sizes are $q = 2^k$
  with $k$ odd, since that is exactly when the self-dualization applies.
\end{enumerate}

So the honest summary of what this blueprint assumes, beyond Mathlib, is one classical
algebra lemma and one rigidity theorem used through a single consequence. That is a
smaller surface than the paper's, and it is smaller because of the value form and direct
repetition rather than by accident.
\end{remark}

\begin{theorem}[Explicit succinct Cook--Levin]
\label{thm:succinct-sat}
\ledgernode{1.3.1}
\uses{def:succinct,def:decider}
There is a polynomial-time algorithm that, given a decider $\decider$, an index $n$,
a time bound $T$ (in binary) and strings $x, y$, outputs a succinct description (in the
sense of Definition~\ref{def:succinct}, via Boolean formulas) of a 3SAT instance
$\varphi$ of size $2^{\poly(\log T)}$ such that $\varphi$ is satisfiable if and only if
there exist strings $a, b$ of length at most $T$ with $\decider(n, x, y, a, b) = 1$
within $T$ steps; moreover satisfying assignments of $\varphi$ encode the
computation tableau together with $(a,b)$.
\end{theorem}

\paragraph{Source.} Cook~\cite{Coo71} and Levin for the tableau reduction; Papadimitriou
and Yannakakis~\cite{PY86} for succinct representations and their complexity;
the explicit form needed here is~\cite[Section 10]{JNVWY20}.

\paragraph{Comments.} This is the classical core of answer reduction: the players prove
satisfiability of $\varphi$ (arithmetized over a low-degree extension) instead of
sending $(a,b)$. The statement is long but entirely combinatorial. Cook--Levin has been
formalized in other proof assistants; a Lean formalization with the explicit-uniformity
guarantees above is a substantial but well-specified sub-project.

A caution on the construction, from the verification campaign~\cite{Audit26} and not
checked here: in the Tseitin encoding of~\cite{NW19}, which the explicit form
of~\cite[Section 10]{JNVWY20} follows, the formula is reported to omit the output-wire
conjunct, so that it is satisfiable for every input and the central equivalence of the
construction fails as written. The statement above is end-to-end and so does not inherit
the defect textually, but a formalization of it must supply what the repair needs: the
output-wire conjunct, the out-degree-zero hypothesis that yields individual degree $6$,
and the exact variable count on which the later PCP parameters depend --- none of which
the statement currently fixes.

\subsection{Self-dual normal bases}
\label{sec:rr-self-dual-basis}

\effortMedium

\tnote{Is this needed?}

\begin{lemma}[Deterministic self-dual normal basis]
\label{lem:self-dual-basis}
\ledgernode{1.1.6, 1.1.6.1}
There is a deterministic algorithm that, given an odd integer $k > 0$, outputs in time
$\poly(k)$ a self-dual normal basis of $\F_{2^k}$ over $\F_2$ together with its
multiplication tables.
\end{lemma}

\paragraph{Source.} Combines Shoup's deterministic irreducible-polynomial
construction~\cite{Sho90}, Lenstra's isomorphism algorithm~\cite{Len91}, and the
classical existence and construction of self-dual normal bases in characteristic $2$
for odd $k$; see the discussion in~\cite[Section 3]{JNVWY20}.

\paragraph{Comments.} Needed so that verifiers can do field arithmetic over
$\F_{2^k}$ using only $\F_2$-data (the ``downsize'' maps), uniformly and
deterministically. This is why field sizes are restricted to $q = 2^k$, $k$ odd
(``admissible''). Self-contained algebra/algorithmics; independent of everything
quantum.

\subsection{The Schwartz--Zippel lemma}
\label{sec:rr-schwartz-zippel}

\effortMathlib

\begin{lemma}[Schwartz--Zippel]
\label{lem:schwartz-zippel}
Let $f \neq g : \F_q^m \to \F_q$ be polynomials of total degree at most $d$. Then
$\Pr_{x \sim \F_q^m}[ f(x) = g(x) ] \leq d/q$.
\end{lemma}

\paragraph{Source.} \cite{Sch80,Zip79}. A version is available in Mathlib
(multivariate polynomials over finite fields).

\paragraph{Comments.} Used throughout the low-degree machinery (distance of the
Reed--Muller code, soundness of arithmetization). The variant for individual degree
should be derived alongside it.

\subsection{Approximating the value from below ($\MIPstar \subseteq \RE$)}
\label{sec:rr-value-approx}

\effortMedium

\begin{lemma}[The value is approximable from below]
\label{lem:value-lower-approx}
\ledgernode{1.1.7.2, 1.1.7.2.2, 1.1.7.2.5, 1.1.7.2.6}
\uses{def:game-description,def:sync-value,def:q-value,def:re,lem:norm-two-psd,lem:perturbation-split}
The set $\{ (g, p) : g \text{ a game description}, p \in \mathbb{Q},
\synval(\game_g) > p \}$ is recursively enumerable, and likewise for $\valstar$.
\end{lemma}

\paragraph{Source.} Folklore; see~\cite[Section 12]{JNVWY20}. Enumerate
tensor-product strategies with entries in $\mathbb{Q}[i]$ and evaluate; density of
such strategies among all tensor-product strategies gives completeness of the
enumeration. That argument gives the $\valstar$ half; the $\synval$ half needs its own
enumeration and density argument, since rounding the entries of a synchronous strategy
does not preserve projectivity or synchronicity.

\begin{lemma}[Decidable candidate set: a norm constraint is two psd tests]
\label{lem:norm-two-psd}
\lean{MIPRE.ValueApprox.posSemidef_realSmul_one_add_and_sub_iff}
\leanok
\ledgernode{1.1.7.2.1}
Let $M$ be a Hermitian matrix and $c$ a real number. Then $c\Id + M \succeq 0$ and
$c\Id - M \succeq 0$ both hold if and only if $-c\norm{x}^2 \leq \bra{x} M \ket{x}
\leq c \norm{x}^2$ for every vector $x$ --- which for Hermitian $M$ says
$\norm{M} \leq c$.
\end{lemma}

\begin{proof}
\leanok
Both directions are the defining property of positive semidefiniteness applied to the
quadratic form, read once forwards and once backwards; the scalar step is
$\bra{x}(c\Id \pm M)\ket{x} = c\norm{x}^2 \pm \bra{x} M \ket{x}$.
\end{proof}

\begin{lemma}[The perturbation split]
\label{lem:perturbation-split}
\uses{def:q-value}
\lean{MIPRE.ValueApprox.dotProduct_mulVec_perturb,
MIPRE.ValueApprox.kronecker_sub_kronecker,
MIPRE.ValueApprox.dotProduct_kronecker_perturb}
\leanok
\ledgernode{1.1.7.2.3}
\ledgernode{1.1.7.2.4}
For matrices $M, N$ and vectors $\psi, \phi$,
\[
\bra{\psi} M \ket{\psi} - \bra{\phi} N \ket{\phi}
  = \bra{\psi} (M - N) \ket{\psi} + \bra{\psi - \phi} N \ket{\psi}
    + \bra{\phi} N \ket{\psi - \phi},
\]
and $A \ot B - A' \ot B' = (A - A') \ot B + A' \ot (B - B')$; combining the two splits a
bipartite Born expectation in which both players' operators and the state have moved into
four terms, each moving exactly one of $A$, $B$ and $\psi$.
\end{lemma}

\begin{proof}
\leanok
Expand and cancel. Both identities are exact: no norm and no estimate occurs in either,
and the analytic bounds are obtained afterwards by bounding the terms separately.
\end{proof}

\paragraph{Comments.} This is the easy inclusion $\MIPstar \subseteq \RE$ and one half
of Theorem~\ref{thm:mipstar-eq-re}. The pipeline consumes only the $\valstar$ half, as the
semidecision procedure of Remark~\ref{rem:compression-abstract}. The commuting-operator analogue (approximating
$\valco$ from \emph{above} via the NPA hierarchy~\cite{NPA08}) is \emph{not} needed for
the main theorem, but becomes relevant for the downstream consequences in
Section~\ref{ch:downstream}.
